\documentclass[10pt,a4paper]{amsart}
\usepackage[margin=19mm]{geometry}
\usepackage{amsmath,amssymb,mathtools}
\usepackage[scr=rsfs,cal=euler]{mathalfa}
\usepackage{xfrac}
\usepackage[unicode,hidelinks,bookmarks=false]{hyperref}
\newcommand{\ie}{i.e.\ }
\newcommand{\cf}{cf.\ }
\newcommand{\resp}{respectively\ }
\newcommand{\Subsection}{\subsection}
\providecommand{\nobreakdash}{\nobreak\hbox{-}}
\setlength{\emergencystretch}{2em}
\multlinegap0pt
\usepackage{bm}
\usepackage{bbm}
\usepackage{dsfont}
\usepackage{xspace}
\usepackage{cancel}
\usepackage{mathtools}
\usepackage{cleveref}
\usepackage{amsthm}
\makeatletter
\@ifundefined{theorem}{\newtheorem{theorem}{Theorem}}{}
\@ifundefined{corollary}{\newtheorem{corollary}[theorem]{Corollary}}{}
\@ifundefined{lemma}{\newtheorem{lemma}[theorem]{Lemma}}{}
\makeatother
\makeatletter
\expandafter\ifx\csname align\endcsname\relax
\renewenvironment{align}{\linenomathNonumbers\oldalign}{\oldendalign\endlinenomath}
\fi
\expandafter\ifx\csname equation\endcsname\relax
\renewenvironment{equation}{\linenomathNonumbers\oldequation}{\oldendequation\endlinenomath}
\fi
\makeatother
\title[Remarks on diagonal dimension for algebraic stacks]{Remarks on diagonal dimension for algebraic stacks}
\author{Pat Lank, Fei Peng}
\date{Source: arXiv:2605.13416v2 (CC BY 4.0)}
\begin{document}
\maketitle

\noindent\emph{Source and licence.} Remarks on diagonal dimension for algebraic stacks. Version arXiv:2605.13416v2 (CC BY 4.0), distributed under the Creative Commons Attribution 4.0 International licence, as recorded in the arXiv licence metadata for this exact version and shown on the abstract page https://arxiv.org/abs/2605.13416v2. Full rights notice: licenses/arxiv-2605.13416-CC-BY-4.0.txt.\par\smallskip

\noindent\textit{Editorial scope.} Counted unit: the Corollary whose source label is \texttt{cor:derived\_resolution} in the article (source0.tex lines 501--512, source label \texttt{cor:derived\_resolution}), delivered together with the article's own complete original proof of it (source0.tex lines 514--536, one \texttt{proof} environment of 23 lines). No mathematics is written, repaired or re-derived: statement and proof are the article's own text line for line. Required context retained uncounted: lem:derived\_resolution\_of\_the\_diagonal (lines 472--483). Every definition, display and label of those lines is retained unchanged. Omitted from this package and not redistributed: 1--471, 484--500, 513--513, 537--786. Nothing inside a retained excerpt is omitted or altered: every retained body is delivered exactly as the article's lines are written, so no omission receipt of any kind accompanies this package. Citations: the retained excerpts carry no citation command at all, so no bibliography entry of the article is transcribed and no reference is redistributed. Reference adaptation: none was needed and none was made; every reference inside the delivered unit resolves inside it, so no unresolved reference or citation is produced. Printed slips kept literal: no printed word, symbol, hypothesis or number of the counted statement or of its proof was altered. Layout and class: the article's own class is \texttt{amsart} with packages including amssymb, amsthm, enumitem, colonequals, tikz-cd, microtype, mathtools, ulem, Baskervaldx, newtxmath, mathalfa, geometry, xcolor, hyperref; the wrapper is the lane's pinned \texttt{amsart} editorial wrapper at 10pt on A4, which restates the theorem-like environments the retained text needs and adds the article's own macro definitions that the retained text needs (none), with extra wrapper packages (bm, bbm, dsfont, xspace, cancel, mathtools, cleveref). The delivered numbering is the unit's own. Assets: no retained span contains a figure environment, a graphics-inclusion command, a PDF-page inclusion, a table or a TikZ picture; no image file is copied into this package, the package declares no asset dependency and no image file is redistributed with it. Licence: version arXiv:2605.13416v2 of this article is distributed under the Creative Commons Attribution 4.0 International licence, as recorded in the arXiv licence metadata for this exact version and shown on the abstract page https://arxiv.org/abs/2605.13416v2. A full rights notice is delivered at licenses/arxiv-2605.13416-CC-BY-4.0.txt. Characters: every retained excerpt is pure ASCII, so the delivered engine, XeLaTeX with its default Latin Modern text font, reports no missing character.
\begin{lemma}
    \label{lem:derived_resolution_of_the_diagonal}
    Let $R$ be a regular ring of finite Krull dimension and $\mathcal{X}$ be a Noetherian algebraic stack with quasi-finite and separated diagonal over $R$. Let $K$ be a complex in $D_{\operatorname{coh}}^b(\mathcal{X}\times_R\mathcal{X})$. Let $g_1,g_2\colon \mathcal{X}\times_R\mathcal{X}\to\mathcal{X}$ be the canonical projections. If $\mathcal{X}$ is smooth and finitely presented over $R$, then there exists a distinguished triangle in $D_{\operatorname{coh}}^b(\mathcal{X}\times_R\mathcal{X})$,
    \begin{displaymath}
        B_1\longrightarrow K_0\longrightarrow K\longrightarrow B_1[1]
    \end{displaymath}
    such that $K_0\to K$ induces surjective maps on the cohomology sheaves and
    \begin{displaymath}
        K_0=\bigoplus_{j\in J}\mathbf{L}g_1^\ast F_j\otimes^{\mathbf{L}}\mathbf{L}g_2^\ast G_j[n_j]
    \end{displaymath}
    for some finite set $J$ and coherent sheaves $F_j$ and $G_j$ on $\mathcal{X}$ and integer $n_j$ for every $j\in J$. Moreover, if $K\in D_{\operatorname{coh}}^b(\mathcal{X}\times_R\mathcal{X})^{\leq 0}$, then so is $B_1$. (Here, $D_{\operatorname{coh}}^b(\mathcal{X}\times_R\mathcal{X})^{\leq 0}$ are those complexes $E$ with bounded coherent cohomology and that satisfy $\mathcal{H}^n (E) \cong 0$ for $n>0$).
\end{lemma}

\begin{corollary}
    \label{cor:derived_resolution}
    In the situation of \Cref{lem:derived_resolution_of_the_diagonal}, there exists a distinguished triangle in $D_{\operatorname{coh}}^b(\mathcal{X}\times_R\mathcal{X})$,
    \begin{displaymath}
        C_n\longrightarrow K\longrightarrow B_{n+1}[n+1]\longrightarrow C_n[1]
    \end{displaymath}
    for every $n\geq 0$ such that $C_n\in\langle\oplus_i^{n} K_i \rangle_{n+1}$ where
    \begin{displaymath}
        K_i=\bigoplus_{j\in J_{i}}\bigg(\mathbf{L}g_1^\ast F_i^j\otimes^{\mathbf{L}}\mathbf{L}g_2^\ast G_i^j[n_j]\bigg)
    \end{displaymath}
    for some finite sets $J_i$ and coherent sheaves $F_i^j$ and $G_i^j$ on $\mathcal{X}$ for every $j\in J_i$ and for every $0\leq i\leq n$. If $K\in D_{\operatorname{coh}}^b(\mathcal{X}\times_R\mathcal{X})^{\leq 0}$, it can be arranged that $B_i\in D_{\operatorname{coh}}^b(\mathcal{X}\times_R\mathcal{X})^{\leq 0}$ for every $0\leq i\leq n+1$.
\end{corollary}

\begin{proof}
    We proceed by induction $n$. If $n=0$, we have a distinguished triangle
    \begin{displaymath}
        B_1\longrightarrow K_0\longrightarrow K=B_0\longrightarrow B_1[1]
    \end{displaymath}
    by \Cref{lem:derived_resolution_of_the_diagonal}. Moreover, if $K=B_0\in D_{\operatorname{coh}}^b(\mathcal{X}\times_R\mathcal{X})^{\leq 0}$, then it can be arranged that $B_1$ is so. In this case, we have $C_0\simeq K_0$, and the statement is trivially true. Suppose there is a distinguished triangle
    \begin{displaymath}
        C_k\longrightarrow K\longrightarrow B_{k+1}[k+1]\longrightarrow C_k[1]
    \end{displaymath}
    for some $k$ such that $C_k\in\langle\oplus_i^{k} K_i \rangle_{k+1}$ where the $K_i$'s are of the desired form. By \Cref{lem:derived_resolution_of_the_diagonal}, we have a distinguished triangle
    \begin{displaymath}
        B_{k+2}\longrightarrow K_{k+1}\longrightarrow B_{k+1}\longrightarrow B_{k+2}[1]
    \end{displaymath}
    where $K_{k+1}$ is of the desired form. Moreover, if $B_{k+1}\in D_{\operatorname{coh}}^b(\mathcal{X}\times_R\mathcal{X})^{\leq 0}$, it can be arranged that $B_{k+2}$ is so. Let $C_{k+1}$ be the cocone of the composition
    \begin{displaymath}
        K\to B_1[1]\to B_2[2]\to\ldots\to B_{k+2}[k+2].
    \end{displaymath}
    By the octahedron axiom, there is a distinguished triangle
    \begin{displaymath}
        K_{k+1}[k]\longrightarrow C_{k}\longrightarrow C_{k+1}\longrightarrow K_{k+1}[k+1].
    \end{displaymath}
    Since $C_k\in\langle\oplus_i^{k+1} K_i\rangle_{k+1}$ and $K_{k+1}[k+1]\in\langle\oplus_i^{k+1} K_i\rangle_{1}$, we have $C_{k+1}\in\langle\oplus_i^{k+1} K_i\rangle_{k+2}$. By induction, the statement follows.
\end{proof}

\end{document}
